\begin{longtable}{p{2.6cm} p{4.6cm} p{2.0cm} p{5.0cm}}
\caption{Métricas previstas para validación.}
\label{tab:metricas-validacion} \\
\toprule
\textbf{Área} & \textbf{Métrica} & \textbf{Unidad} & \textbf{Método de cálculo} \\
\midrule
\endfirsthead

\multicolumn{4}{c}{\tablename\ \thetable{} -- continuación} \\
\toprule
\textbf{Área} & \textbf{Métrica} & \textbf{Unidad} & \textbf{Método de cálculo} \\
\midrule
\endhead

\midrule
\multicolumn{4}{r}{Continúa en la siguiente página} \\
\endfoot

\bottomrule
\endlastfoot

Control & Error estacionario, sobreimpulso y tiempo de establecimiento & \%, °C &
\begin{minipage}[t]{5.0cm}
\vspace{2pt}
$\displaystyle e_{ss} = \dfrac{|SP - T_{estable}|}{SP}\times 100$

\vspace{6pt}
$\displaystyle M_p = \dfrac{T_{max}-SP}{SP}\times 100$

\vspace{6pt}
Calculados a partir de la curva $T(t)$ exportada del PLC/HMI.
\end{minipage} \\[4pt]
Comunicaciones & Latencia, pérdida de mensajes y tiempo de reconexión & ms, \%, s & Captura de tráfico de red con marcas de tiempo; conteo de paquetes esperados frente a recibidos \\
Dosificación & Error absoluto y relativo del volumen dosificado & mL, \% & Comparación entre el volumen programado (por tiempo de activación) y el volumen medido en probeta, en 3 repeticiones \\
Recuperación & Porcentaje de recuperación de agua & \% & Relación entre el volumen recuperado en TK-103 y el volumen total suministrado al ciclo \\
Instrumentación & Error de medición de nivel (LIT-101) & \% & Comparación entre la lectura del sensor y la medición manual de referencia en TK-102, en al menos 3 niveles \\
Agitación & Tiempo de homogeneización de la mezcla (M-101/AG-101) & min & Cronometraje desde el inicio de la agitación hasta lograr una mezcla visualmente homogénea, en 3 repeticiones \\
Seguridad & Tiempo de respuesta del paro de emergencia y de las alarmas & s & Cronometraje o captura de eventos del PLC desde la activación hasta el corte o la señal de alarma \\

\end{longtable}